% Part: normal-modal-logic
% Chapter: axioms-systems
% Section: duals

\documentclass[../../../include/open-logic-section]{subfiles}

\begin{document}

\olfileid{nml}{prf}{dua}

\olsection{دوه‌ګوني \usetoken{P}{formula}}

\begin{defn}\ollabel{def:duals}
  د \Ax{T}، \Ax{B}، \Ax{4} او \Ax{5} له !!{formula}s څخه هر يو
  \emph{دوه‌ګونے} لري؛ لکه لاندې چې ښکاري، هغه د لاندېنۍ لوزې په نښه
  ښودل کېږي:
  \begin{align}
    \tag{\Ax{T_\Diamond}} p & \lif \Diamond p\\
    \tag{\Ax{B_\Diamond}} \Diamond\Box p & \lif p\\
    \tag{\Ax{4_\Diamond}} \Diamond\Diamond p & \lif \Diamond p\\
    \tag{\Ax{5_\Diamond}} \Diamond\Box p & \lif \Box p
    \end{align}
\end{defn}

له پورتنيو دوه‌ګونو !!{formula}s څخه هر يو له اړوند !!{formula}
څخه داسې لاس ته راځي: لومړے د $p$ پر ځاے $\lnot p$ کېږدو، بيا
مقابل نقيض واخلو، وروسته $\lnot\Box\lnot$ په $\Diamond$ او
$\lnot\Diamond\lnot$ په~$\Box$ بدل کړو۔ \Ax{D}، يعنې
$\Box!A \lif \Diamond!A$، په همدې معنا خپله دوه‌ګونی دے۔

\begin{prop}\ollabel{prop:dualsys}
  په \olref[nml][prf][dua]{def:duals} کښې د هر !!{formula}~$!A$ دپاره:
  $\Log{K}!A = \Log{K}!A_{\Diamond}$.
\end{prop}
\begin{proof}
  تمرین۔
\end{proof}
\begin{prob}
  \olref[nml][prf][dua]{prop:dualsys} ثابته کړئ۔
\end{prob}

\end{document}
